% Review
\subsection{Sololearn}

% Review
Sololearn je jednou z nejpopulárnějších aplikací na učení se softwarového inženýrství a programování. Nabízí přes 40 kurzů na různá témata napříč softwarovým inženýrstvím a má napříč platformami přes 58 milionů registrovaných uživatelů.

% Review
Sololearn je k dispozici jako webová aplikace a dále také jako mobilní aplikace na platformách Android a iOS. Pro analýzu použiji mobilní verzi aplikace pro platformu Android.

% Review
\subsubsection{Vzhled a uživatelská zkušenost}

% Odstavec: Obecný dojem (komplexita/přehlednost UI, paleta barev, typografie)

% Review
Design aplikace může na první pohled působit starším dojmem, zejména kvůli použití některých komponentů knihovny Material UI bez úpravy vzhledu a kvůli nedodržení některých doporučených standardů designu. % Write přidat referenci na Mui text density https://m2.material.io/design/layout/applying-density.html#layout
Uživatelské rozhraní místy působí příliš nepřehledně. Co se palety barev týče, používá Sololearn světle modrou jako primární barvu a světle zelenou jako sekundární barvu. Pro pozadí a plochy pak používá bílou a modrošedou barvu. V rámci typografie je napříč celou aplikací použito písmo Fira Sans, jen pro bloky kódu je použito monospace písmo Fira Mono.


% Odstavec: Rozložení prvků na obrazovce (lekce v kurzu, učení)
% Review
Rozložení prvků na obrazovce je poměrně přehledné a logické. Místy může působit nepřehledně kvůli malému prostoru mezi jednotlivými prvky, malé velikosti textu a většímu nahuštění prvků uživatelského rozhraní na jednotlivých stránkách obrazovce.

% Review
Na hlavní stránce kurzu jsou lekce seskupeny do sekcí, které může uživatel kliknutím rozbalit. Díky tomu si může jednoduše získat přehled o celém kurzu a rozbalit části, které ho zajímají.

% Odstavec: Navigace v aplikaci

% Review
Pro navigaci je v aplikaci použita dolní navigační lišta s pěti prvky a také boční navigační lišta s ostatními funkcemi. 

% Odstavec: Srovnání webového a mobilního prostředí

% Review
Webová verze aplikace Sololearn obsahuje méně funkcí, než verze mobilní. V mobilní aplikaci je kromě základních funkcí také možné plnit programátorské výzvy, soutěžit ve výzvách proti ostatním uživatelům, číst články a procházet lekce vytvořené komunitou.  


% Odstavec: Animace a efekty a mikro interakce
% Review
Sololearn využívá animace a efekty velice zřídka. Základní animace jsou použity zejména pro zobrazování prvků na obrazovce. Nejvíce animací je využito v animovaných obrázcích, které jsou zobrazovány po dokončení lekce nebo v úvodním registračním formuláři.

% Review
\subsubsection{Způsoby učení}

% Struktura kurzu (lineární, tematické moduly, career paths, kapitoly atd.)
% Review
Kurzy jsou v aplikaci rozděleny do sekcí, ve kterých se nacházejí jednotlivé lekce a kvízy. Díky této jednoduché struktuře je velice snadné se v kurzech orientovat. Potenciální nevýhodou této jednoduché struktury může být to, že se v delších kurzech nachází příliš mnoho sekcí a uživatel, což může být pro uživatele nepřehledné.  

% Jak se spustí učení a jak vypadá study session
% Review
Pro spuštění učení uživatel přejde na stránku kurzu, kde se u nadcházející lekce nachází tlačítko pro spuštění učení. Samotné učení se velice podobá učení v aplikaci Mimo. Skládá se ze série stránek vysvětlující danou látku a krátkých cvičení, při kterých si uživatel látku procvičí. Po zodpovězení otázky si uživatel může také zobrazit diskuzi a komentáře ostatních uživatelů k danému cvičení. Na rozdíl od aplikace Mimo se po dokončení cvičení nezobrazuje automaticky výstup daného kódu. Uživateli je jen zobrazeno, zda odpověděl správně nebo špatně a může si nechat jeho odpověď vysvětlit pomocí umělé inteligence. Na konci každé lekce je uživateli zobrazeno shrnutí s klíčovými informacemi, které by si měl z lekce odnést.

% Druhy cvičení
% Review
Sololearn nabízí v rámci učení různé druhy cvičení. Nejčastěji se jedná o doplňování slov z nabídky do kódu nebo textu. Dalším častým typem cvičení je výběr správné odpovědi z nabídky. Sololearn také nabízí cvičení, v rámci kterého může uživatel psát celý kód a zobrazovat si jeho výstupy. Celou strukturou učení se tak velice podobá analyzované aplikaci Mimo.

% Review
\subsubsection{Způsoby motivace uživatele}

% Jaké formy gamifikace aplikace používá
% Review
Sololearn pro motivaci uživatelů využívá několik forem gamifikace. Podobně jako v aplikaci Mimo využívá streak uživatele, srdíčka při učení, body XP, ligu učení a virtuální měnu.

% Review
Kromě těchto forem gamifikace jsou uživateli také zobrazovány na různých místech v aplikaci indikátory pokroku, které zobrazují, jakou část jednotlivých kurzů uživatel dokončil.

% Review
\subsubsection{Způsoby monetizace}

% Review
Aplikace nabízí uživatelům dva typy předplatných Sololearn Pro a Sololearn Max.

% Free tier
% Review
Uživatelé bez předplatného mohou procházet standardní lekce ve všech kurzech. Jediným omezením je počet srdíček, které uživatelé mají. Uživatelům jsou také po dokončení lekcí zobrazovány reklamy.

% Pro tier
% Review
Uživatelé s předplatným Pro mají přístup k speciálním lekcím v kurzech, ve kterých si mohou procvičit programátorské úlohy zaměřené na reálné problémy z praxe. Dále mají k dispozici neomezený počet srdíček, nejsou jim zobrazovány reklamy a mají přístup k dodatečným materiálům a kvízům.

% Max tier
% Review
Předplatné Max kromě zmíněných výhod nabízí také možnost uživateli využít umělou inteligenci, aby mu pomohla při učení. Díky ní si uživatel může nechat vysvětlit daný kus kódu, proč byla jeho odpověď špatná, a jak by ji měl opravit. 

% Review
\subsubsection{Silné a slabé stránky}

% Silné stránky
% Review
Silnou stránkou této aplikace je přehlednost a jednoduchá struktura jednotlivých kurzů. Pro obsáhlejší kurzy může být tato struktura nepřehledná, nicméně pro malé a středně velké kurzy, které se v aplikaci vyskytují, může být tato struktura vhodná.

% Review
Další silnou stránkou je samotné učení. Stránky s vysvětlením látky i cvičení se snadno používají. Látka tak může být pro uživatele mnohem stravitelnější, než například u cvičení z konkurenční aplikace Codecademy.

% Review
Poslední silnou stránkou Sololearn je motivace uživatele. Aplikace využívá řadu mechanismů pro motivaci uživatele, které mohou přimět uživatele pravidelně aplikaci používat 

% Slabé stránky
% Review
Jednou ze slabých stránek Sololearn je design některých stránek. Kvůli nedodržení některých standardů designu mohou stránky působit na uživatele příliš nepřehledně a zastaralým dojmem.

% Review
Druhou slabou stránkou aplikace je délka úvodního dotazníku při registraci uživatele. Dotazník obsahuje přes 10 stránek, což může být pro některé uživatele potenciálně odrazující. Založení účtu tak trvá výrazně déle, než u ostatních aplikací. 

% Review
\subsubsection{Další poznatky}

% Zajímavé funkce

% Review
Při učení aplikace na mobilu umožňuje uživateli přejet prstem po obrazovce doleva nebo doprava a tím přecházet mezi již dokončenými cvičeními v rámci dané lekce. Díky tomu se může uživatel velice jednoduše vrátit k předchozím cvičením a stránkám, které vysvětlují danou látku.

% Review
Další funkcí, která stojí za pozornost, je dostupné a snadné přepínání mezi zapsanými kurzy. Na rozdíl od ostatních aplikací, kde uživatel musí pro přepnutí mezi kurzy přejít na samostatnou stránku, může v Sololearn uživatel jednoduše přepnout kurz kliknutím na název kurzu v horní navigační liště, čímž se rozbalí seznam zapsaných kurzů, ze kterých si může uživatel jeden vybrat a přesunout se tak na jeho hlavní stránku. Díky tomu může uživatel snadno přecházet mezi zapsanými kurzy, aniž by se musel přesouvat na samostatnou stránku se seznamem kurzů. To může být velice užitečné zejména pro uživatele, kteří se chtějí učit více kurzů zároveň.  
